\chapter{User variables and equations}\label{chap:user}

\manref{cge_variable} and \manref{cge_equation} add variables and equations to
the calibrated model, written in a small language; \manref{cge_set} and
\manref{cge_parameter} add the sets and parameters they may use (from
version 0.7.1), conditions restrict sums, terms and domains (from version
0.7.2), and \manref{cge_macro} adds macros (from version 0.9.7). From version
0.13.0 parameters may be calibrated by assignments, and a whole model may be
written in the language, without a SAM or the built-in equations
(section~\ref{sec:standalone}).

\section{The language}

An expression is made of
\begin{itemize}
\item numbers, the operators \cmd{+ - * / ^}, parentheses, and the functions
\cmd{exp()}, \cmd{ln()} (or \cmd{log()}) and \cmd{sqrt()}, and from version
0.9.7 \cmd{max()}, \cmd{min()}, \cmd{abs()}, \cmd{sqr()}, \cmd{log10()} and
\cmd{ifthen()} (section~\ref{sec:userfunctions});
\item model and user variables, \cmd{VAR}, \cmd{VAR(e)} or \cmd{VAR(e1, e2)},
at an element or at the index of the equation or of a sum;
\item parameters that \manref{cge_shock} changes, \cmd{par(e)}, and user
parameters, \cmd{upar}, \cmd{upar(e)} or \cmd{upar(e1, e2)}, at their current
values;
\item \cmd{base(VAR(e))}, the base value of a variable;
\item \cmd{sum(i, expression)}, a sum over the set of index \cmd{i}: \cmd{a}
activities, \cmd{c} commodities, \cmd{f} factors, \cmd{h} households, \cmd{d}
households and enterprises, \cmd{i} institutions that receive factor income,
\cmd{n} nests, or a user set, by its name; digits may follow, as in \cmd{c2}
or \cmd{goods2};
\item \cmd{sum((i,j), expression)}, a sum over every pair of elements of two
sets, and \cmd{sum(map(i,j), expression)}, a sum over the pairs of a set of
pairs (from version 0.9.1);
\item \cmd{prod()}, \cmd{smax()} and \cmd{smin()}, the product, the largest
and the smallest value over a set, with the domains of \cmd{sum()}, and
\cmd{card(set)}, the number of its elements (from version 0.9.7);
\item macros, \cmd{name(x, y)}, defined by \manref{cge_macro} (from version
0.9.7);
\item \cmd{ord(i)}, the position of the element of index \cmd{i} in its set,
from 1;
\item \cmd{lag(VAR)} and \cmd{lag(VAR(e), k)}, the value of a variable in the
period before or $k$ periods before, up to four (from version 0.9.2;
section~\ref{sec:userlags});
\item conditions, \cmd{$(condition)}, on sums, terms and domains
(section~\ref{sec:userconditions}).
\end{itemize}

\section{User sets and parameters}

A user set (\manref{cge_set}) is an ordered list of element labels, such as
\cmd{cge_set goods = c, except(com_ser com_hlt)}. Its name is an index like
the letter of a model set: \cmd{cge_variable PREL(services) = PQ(services) /
CPI} adds a variable for each element of \cmd{services}, and
\cmd{sum(goods, PQ(goods) * QG(goods))} sums over the elements of
\cmd{goods}. A user set adds no variables or equations of its own.

A user parameter (\manref{cge_parameter}) has no index, one or two, over
model sets or user sets, and values that the user gives: one value for all
elements, values by element, or a variable of the data in memory. Its values
are part of the model's parameters: every simulation stores them, and
\manref{cge_shock} changes them, as in \cmd{cge_shock hgoal, value(1.2)}, so
scenarios can be written in terms of the user's own parameters. A
parameter's element may also be an instrument of \manref{cge_optimize}.
User sets and parameters are saved with the model, and kept when it is
calibrated again, the parameters at their base values.

From version 0.13.0 a parameter's value may be an expression, as in GAMS's
assignments: \cmd{cge_parameter share(c,a) = qint0(c,a) / sum(c2,
qint0(c2,a))} evaluates the expression for each element, with the indexes
written out, and \cmd{cge_parameter tmr(c)$(tm(c) > 0) = tm(c) / 2} only for
the elements that meet the condition, the others taking the default. The
expression may use numbers, user parameters, the model's parameters, the
base values of variables, \cmd{base(QM(c))}, sums, conditions and the
functions of the language, but not variables themselves. A parameter that
exists may be assigned again, over the elements of a condition, until a
variable, equation or macro uses it. The values are computed once, when the
parameter is assigned: a parameter computed from another does not follow
later shocks to it, which belong in the equations. Data in wide form, as a
SAM in memory, give a parameter with two indexes,
\cmd{cge_parameter sam(ac,ac2), from(act_agr-row) labels(account)}, and a
string variable gives the elements of a set, \cmd{cge_set ac, from(account)}.

\section{Conditions}\label{sec:userconditions}

From version 0.7.2, a condition restricts a sum, a term or the domain of a user
variable or equation, with the dollar operator of GAMS:
\begin{itemize}
\item \cmd{sum(c$(tm(c) > 0), PM(c) * QM(c))} sums over the commodities with
an import tariff;
\item \cmd{(PQ(c) * QG(c))$(goods(c))} is the term where \cmd{c} belongs
to the user set \cmd{goods}, and 0 elsewhere;
\item a name with a condition, as in \cmd{cge_variable TMV(c)$(tm(c) > 0) =}
\cmd{tm(c) * PM(c) * QM(c)}, defines a user variable or equation for the
elements that meet the condition only.
\end{itemize}
A condition compares values with \cmd{> < >= <= ==} and \cmd{!=}, and combines
comparisons with \cmd{&}, \cmd{|} and \cmd{!}; it may test whether an element
belongs to a set, \cmd{goods(c)}, whether the index stands for a given element,
\cmd{sameas(c, com_hlt)}, or the position of the element, \cmd{ord(c)}.
Arithmetic is evaluated first, then comparisons, \cmd{!}, \cmd{&} and
\cmd{|}.

A condition is evaluated when the variable or equation is added, and again
when the model is calibrated, with the variables and parameters at their base
values. It decides which terms and equation elements the model contains; a
shock to a parameter that a condition uses changes the value of the terms that
remain, not which terms remain. The system therefore has the same size in every
simulation, dynamic run and sensitivity draw, and the GAMS program writes out
the same terms. GAMS itself would evaluate the condition again at each solve
statement, with the current values of the parameters; the program that
\manref{cge_gams} writes contains the terms already chosen, so the two agree.
An element of a user variable outside its domain does not exist: it is left
out of the results, it is 0 where an expression refers to it, and GAMS fixes
it at 0. Because \cmd{$}\argu{name} is a global macro on a Stata command
line, the condition is always written in parentheses after \cmd{$}.

\section{Two indexes and sets of pairs}\label{sec:userpairs}

From version 0.9.1, a user variable or equation may have two indexes, as the
model's household demand $\V{QH}_{c,h}$ does: \cmd{cge_variable X(c,h) = PQ(c)
* QH(c, h)} adds a variable for every commodity and household, and a
condition may restrict it to some pairs. A set of pairs, \cmd{cge_set map(a,c)
= act_agr.com_agr act_man.com_man}, records a mapping between two sets, as a
GAMS set over two indexes does. A sum runs over its pairs,
\cmd{sum(map(a,c), PA(a) * QA(a))}, and in a condition, \cmd{map(a,c)}, it
restricts sums, terms and the domains of variables and equations with two
indexes to its pairs. Sums over two indexes and over sets of pairs are written
out one index at a time, so they take conditions as other sums do. In the
results the second index is the column \cmd{element2}; in \manref{cge_shock}
an element is written \cmd{e1.e2}; in the GAMS program the variables are
declared over two sets of accounts, \cmd{(AC, AC)}, with the pairs outside
their domains fixed at 0.

From version 0.25.0 a symbol may have three or more indexes, as GAMS's do:
\cmd{cge_parameter a(i,j,k) = 100 * ord(i) + 10 * ord(j) + ord(k)},
\cmd{cge_variable X(i,j,k)} and \cmd{cge_equation EX(i,j,k): X(i,j,k)^2 =
a(i,j,k)}. A sum runs over any tuple of indexes, \cmd{sum((j,k), X(i,j,k))},
and a set of tuples, \cmd{cge_set m(i,j,k) = p1.q1.r1 p2.q3.r1}, serves in
sums, \cmd{sum(m(i,j,k), a(i,j,k))}, and in conditions, \cmd{m(i,j,k)}, as a
set of pairs does. An element is written \cmd{p2.q3.r1} in
\manref{cge_shock}, \manref{cge_bound} and \cmd{values()}; long data take
one label variable for each index, \cmd{from(v) labels(li lj lk)}. In the
results the second element column holds the later labels joined by periods,
\cmd{q3.r1}, and the GAMS program declares the variables over \cmd{(AC, AC,
AC)}.

Labels may be the names of indexes, as the years \cmd{t1}, \cmd{t2} and
\cmd{t3} of a set whose alias is \cmd{t2}. From version 0.22.1 an equation
over two indexes puts them in correctly, and from version 0.23.0 so does a
sum over an alias within an equation, a defined variable or a parameter's
definition: in \cmd{cge_equation EX(g): X(g) = sum(g2, a(g,g2))}, with
labels \cmd{g1} to \cmd{g3}, the element \cmd{g2} put in for \cmd{g} is
marked, so that the sum over \cmd{g2} does not replace it again. Before,
$X_{g2}$ summed $a_{j,j}$ over $j$. A sum over the equation's own index is
an error.

\section{Lags}\label{sec:userlags}

From version 0.9.2, \cmd{lag(VAR)} is the value of a variable in the period
before, and \cmd{lag(VAR(e), k)} its value $k$ periods before, for $k$ from 1 to
4, as $\V{VAR}_{t-k}$ in a model indexed by time. In a period of a dynamic run
(chapter~\ref{chap:dynamics}) the values are those of the periods solved
before it, and those of the base before the first period. Within a period they
are numbers, as the capital stock is, so a lag adds no variable and its
derivative is zero; in a static simulation the period before is the base. A
stock that the model does not update itself,
\begin{equation*}
S_t = (1 - \delta) S_{t-1} + I_t,
\end{equation*}
is written \cmd{cge_equation s: S = (1 - 0.05) * lag(S) + ...}, with \cmd{S}
a user variable given a value. Each period stores the lagged values it used,
and the GAMS program declares them as parameters, \cmd{LAG1_U_S} for
\cmd{lag(S)}, which each period of a dynamic run sets before it is solved.

From version 0.10.2, \cmd{lead(VAR)} and \cmd{lead(VAR(e), k)} are the values
of later periods, $\V{VAR}_{t+k}$, in a run with perfect foresight
(section~\ref{sec:foresight}), which solves the periods together: across
periods, a lag or a lead is then a variable of another period, with its
derivative. Beyond the last period a lead follows the run's terminal
condition; a recursive run cannot have leads, and in a static simulation they
are the base values. The value of a unit of capital that earns the real
rental rate $R$ and depreciates at rate $\delta$, at interest rate $r$,
\begin{equation*}
Q_t = \frac{R_{t+1} + (1 - \delta)\,Q_{t+1}}{1 + r},
\end{equation*}
is \cmd{cge_equation q: Q = (lead(RK) + 0.95 * lead(Q)) / 1.05}, with
\cmd{RK} defined as \cmd{W(cap) / CPI}. The names \cmd{tperiod}, \cmd{tfirst}
and \cmd{tlast} are the period of a dynamic run and 1 in its first or last
period (0 otherwise, and all three 0 in a static simulation), so that
\cmd{ifthen(tlast, $x$, $y$)} gives the last period an equation of its own;
since they change by period, conditions (\cmd{$}), which hold in every
period, may not use them. The GAMS program declares the values of later
periods as parameters \cmd{LEAD1_U_Q} and the period as the scalars
\cmd{TPERIOD}, \cmd{TFIRST} and \cmd{TLAST}.

\section{Functions and macros}\label{sec:userfunctions}

From version 0.9.7, the language has the functions of GAMS that a square
system solves, as GAMS writes them: \cmd{max(x, y, ...)} and
\cmd{min(x, y, ...)} of two or more arguments, \cmd{abs(x)}, \cmd{sqr(x)},
\cmd{log10(x)} and \cmd{ifthen(condition, x, y)}; the products and extremes
over sets, \cmd{prod()}, \cmd{smax()} and \cmd{smin()}; and \cmd{card()}, the
number of elements of a set. The sums of GAMS have the same domains here:
one index, two indexes, a set of pairs, each with a condition. A product over
no element is 1; a maximum or minimum over none is an error, since GAMS's
$\mp\infty$ is no value for a model.

The condition of \cmd{ifthen()} uses parameters, lags and numbers, which are
given within a simulation; a shock to the parameter switches the branch.
GAMS's CNS and MCP models refuse relations between endogenous variables, which
make the system discontinuous. A regime switch on endogenous variables is an
inequality with its paired variable (section~\ref{sec:userineq}), the
complementarity form in which GAMS solves it too.

The derivatives remain exact. \cmd{max()}, \cmd{min()} and the branch of
\cmd{ifthen()} pass on the derivatives of the argument they take, the first of
equal arguments at a kink; \cmd{abs()} multiplies by the sign of its argument,
\cmd{sqr()} by $2x$ and \cmd{log10()} by $1/(x \ln 10)$. A model whose
solution lies at a kink of \cmd{max()}, \cmd{min()} or \cmd{abs()} may need
more Newton iterations, in Stata as in GAMS.

A macro (\manref{cge_macro}), as GAMS's \cmd{\$macro}, names an expression with
parameters: \cmd{cge_macro ces(x, y, d, r) = (d * x^(-r) + (1 - d) *
y^(-r))^(-1 / r)} makes \cmd{ces(QD(c), QM(c), 0.5, 0.5)} a CES aggregate of
domestic goods and imports. Macros are replaced by their expressions, with the
arguments in parentheses, before the indexes and sums are written out, so
indexes may be arguments and a macro adds nothing to the system. In the GAMS
program the expressions are written out, and the functions keep their GAMS
names, with \cmd{eq}, \cmd{ne}, \cmd{and}, \cmd{or} and \cmd{not} in the
conditions of \cmd{ifthen()}.

\section{Replacing built-in equations}\label{sec:userreplace}

From version 0.9.3, a user equation may take the place of elements of a
built-in equation block of chapter~\ref{chap:equations}, with
\cmd{replace(}\argu{block}\cmd{)}:
\begin{stcode}
cge_equation qe: QE(com_agr) = base(QE(com_agr)), replace(T2(com_agr))
\end{stcode}
holds the exports of agricultural goods at the base, in place of their export
supply. The replaced elements leave the system and the user equation joins it,
so the system stays square and the closure is unchanged. The blocks that may
be replaced, by domain, are
\begin{cgetable}{p{0.14\linewidth}p{0.80\linewidth}}
\toprule
Domain & Blocks\\
\midrule\endhead
\bottomrule\endlastfoot
none & \eqref{eq:H1}, \eqref{eq:I12}, \eqref{eq:I14}, \eqref{eq:S4}\\
$c$ & \eqref{eq:P1} to \eqref{eq:P4}, \eqref{eq:Q6}, \eqref{eq:Q6A}, \eqref{eq:T1} to \eqref{eq:T10}, \eqref{eq:T3E}, \eqref{eq:I6}, \eqref{eq:I7}, \eqref{eq:I13}\\
$a$ & \eqref{eq:P5}, \eqref{eq:P6}, \eqref{eq:Q1}, \eqref{eq:Q2}, \eqref{eq:Q1C}, \eqref{eq:Q2C}, \eqref{eq:Q3}\\
$f$ & \eqref{eq:H2}, \eqref{eq:H3}, \eqref{eq:L1}, \eqref{eq:L2}\\
$h$ & \eqref{eq:I4}\\
$d$ & \eqref{eq:I3}, \eqref{eq:I10}, \eqref{eq:I11}\\
$(f,a)$ & \eqref{eq:Q4}\\
$(n,a)$ & \eqref{eq:N1}, \eqref{eq:N2}\\
$(c,a)$ & \eqref{eq:Q5}, \eqref{eq:N3}\\
$(a,c)$ & \eqref{eq:Q6B}, \eqref{eq:P5A}\\
$(i,f)$ & \eqref{eq:I2}\\
$(c,h)$ & \eqref{eq:I5}\\
$(a,h)$ & \eqref{eq:I5H}\\
\end{cgetable}
Within a block written in variants, such as \eqref{eq:Q1C} for CES and
Cobb--Douglas technologies, the element is replaced whatever its variant. The
price indexes, factor incomes, government revenue and market clearing,
equations \eqref{eq:P7} to \eqref{eq:P9}, \eqref{eq:I1}, \eqref{eq:I8},
\eqref{eq:I9}, \eqref{eq:S1} to \eqref{eq:S3} and \eqref{eq:S5}, have
derivatives coded in the package and cannot be replaced; they can be changed through the closure, or
with user variables and equations that add to them. A replacement changes the
model, and must keep its accounting consistent: the residual of Walras's law,
\cmd{WALRAS}, which every simulation reports, shows when it does not. In the
GAMS program, the parameter \cmd{XRP} marks the replaced elements, which the
equations of \cmd{cge_model.gms} leave out.

\section{Declared nests}\label{sec:usernests}

From version 0.11.0, \manref{cge_nest} declares a tree of CES or CET nests over
variables of the model, as the production blocks of MPSGE do
\citep{rutherford1999}, and writes it as user variables and equations. A cost
block named $B$, with output $Y$, price $P^Y$, efficiency $E$ and leaves $j$,
quantities $x_j$ in units $u_j$ at prices $\pi_j$, is written as
\begin{align*}
w_j &= \pi_j/u_j, \qquad X_j = u_j\,x_j, \qquad
V_j = w^0_j X^0_j \quad (\text{\cmd{B_w}}j,\ \text{\cmd{B_x}}j),\\
P_n &= \Bigl[\sum_{j \in n} \tfrac{V_j}{V_n}\,\bigl(w_j/w^0_j\bigr)^{1-s_n}
       + \sum_{m \in n} \tfrac{V_m}{V_n}\,P_m^{\,1-s_n}\Bigr]^{1/(1-s_n)}
       \quad (\text{\cmd{B_P}}, \text{\cmd{B_}}n\text{\cmd{_P}}),\\
Q_m &= V_m\,R_n\,\bigl(P_n/P_m\bigr)^{s_n} \quad (\text{\cmd{B_}}m\text{\cmd{_Q}}),
\qquad
X_j = X^0_j\,R_n\,\Bigl(\frac{P_n\,w^0_j}{w_j}\Bigr)^{s_n} \quad (\text{\cmd{B_d}}j),\\
P^Y &= P^{Y0}\,P_B/E \quad (\text{\cmd{B_z}}),
\end{align*}
where $V_n$ is the value at base prices of the leaves within nest $n$, $m$ a
nest within $n$, $R_B = (Y/Y^0)/E$ for the block and $R_m = Q_m/V_m$ for a
nest within it. Every $^0$ value is written as \cmd{base()} of a variable or,
from version 0.11.0, of a parameter, so the block holds at the base and is
calibrated again with the model. At $s_n = 1$ the index is a geometric mean
and at 0 a weighted sum; a revenue block has the exponent $1 + t_n$, its
leaves $X_j = X^0_j R_n (w_j/(w^0_j P_n))^{t_n}$, $R_B = (Y/Y^0)\,E$ and
$P^Y = P^{Y0} P_B E$. A leaf exists where its quantity is not zero at the base.
The leaves' equations and $P^Y$ replace the built-in equations named in
\cmd{replace()}, for value added \eqref{eq:Q4} and \eqref{eq:Q3}, for exports
and domestic sales \eqref{eq:T1}, \eqref{eq:T2} and \eqref{eq:P4}. A block of
one nest at the built-in elasticity is the built-in CES or CET function, in
another form; the test suite checks this, and a block of two nests against the
built-in nests of value added (section~\ref{sec:nests}). The tax on factor use,
\cmd{tf(}$f$\cmd{,} $a$\cmd{)}, may be used in expressions from version 0.11.0.

\section{Inequalities and bounds}\label{sec:userineq}

From version 0.9.0, a user equation may be an inequality, $\ell(x) \le r(x)$
or $\ell(x) \ge r(x)$, written with \cmd{=l=} and \cmd{=g=} as in GAMS and paired
with a variable, as in a mixed complementarity problem:
\begin{stcode}
cge_equation quota: QM(com_man) =l= 1.02 * base(QM(com_man)), perp(RM(com_man))
\end{stcode}
The paired variable, fixed under the closure, stays at its level while the
inequality has room to spare and moves when it binds. Regime switches that a
GAMS model writes as complementarity conditions, such as quotas with their
premiums, capacity limits with their rents, floors on spending with an
adjusting instrument, or a minimum price, can so be written in the model's own
terms, with sets, parameters and conditions. Each inequality becomes a
constraint of chapter~\ref{chap:constraints} on its slack variable, and is
solved, reported and exported to GAMS in the same way.

A user variable may be exogenous, \cmd{cge_variable TAU, value(0) fixed}, or
exogenous at a bound from which a paired inequality may move it,
\cmd{cge_variable TAU, lower(0)}: the analogue of a positive variable in
GAMS. A defined variable may be held within bounds,
\cmd{cge_variable PR = PQ(com_agr) / CPI, lower(0.95)}, which makes it the
larger of its expression and the bound.

With variables at a bound and inequalities paired with them, a standalone
model (section~\ref{sec:standalone}) states a mixed complementarity problem
as GAMS's MCP models do: $x \ge 0$, $F(x) \ge 0$, $x'F(x) = 0$ is
\begin{stcode}
cge_variable X(i), lower(0)
cge_equation F(i): f(X) >= 0, perp(X(i))
\end{stcode}
and \manref{cge_bound} gives the start, \cmd{cge_bound X(x1), level(1)}, as
GAMS's \cmd{.l}. From version 0.14.0 the semismooth Newton method of
section~\ref{sec:cons-solution} solves every pair at once.

\section{Variables and equations}

\cmd{cge_variable NAME(i) = expression} adds the variable $U_i$ and the
equation
\begin{equation*}
U_i = e_i(x)
\end{equation*}
for every element $i$ of the set; its base value is $e_i(x^0)$.
\cmd{cge_variable NAME(i), value(v)} adds a free variable without an equation,
and \cmd{cge_equation NAME(i): lhs = rhs} an equation without a variable; each
must be paired with the other, so that the closure stays square.

A defined variable swapped into the closure (\cmd{cge_closure, swap(NAME VAR)})
fixes the expression and frees \cmd{VAR}. For example, real government savings
$\V{GSAV}/\V{CPI}$ held at their base value, with the direct tax rate
adjusting, is under the CPI numeraire the closure \cmd{fixgsav}; the test
suites check that the solutions agree to $10^{-9}$. Nominal government
consumption, $\sum_c \V{PQ}_c\,\V{QG}_c$, fixed with volumes adjusting is
another closure rule that the standard closures do not offer.

\section{Standalone models}\label{sec:standalone}

From version 0.13.0, \cmd{cge_model new} (\manref{cge_model}) starts a model
written whole in the language, as a GAMS program is: sets, parameters,
variables and equations, without a SAM or the built-in equations of
chapter~\ref{chap:equations}. Its sets have labels of their own or from the
data in memory, its parameters are calibrated by assignments, its variables
start at the levels of \opt{}{value()}, expressions of the parameters, and
the exogenous ones are \opt{}{fixed}. The letters of the model's sets are not
indexes in it, so its sets take any names, and the index of a set with a
digit after it, such as \cmd{i2}, runs over the same set, as an alias. The
model's variables are its user variables, in the order they were added, and
its equations two blocks, the user equations and the constraints of their
inequalities:
\begin{equation*}
\ell_k(x) - r_k(x) = 0, \qquad k = 1, \dots, K,
\end{equation*}
one for each element written out. A variable is free unless it was added
\opt{}{fixed}; \cmd{cge_closure, swap()} exchanges a fixed variable for a free
one. \manref{cge_solve} solves the system from the starting levels, which need
not solve it, by Newton's method with the exact Jacobian of
section~\ref{sec:jacobian} and, for large systems, the sparse LU of
section~\ref{sec:sparse}; there is no reduced system, every element being an
equation of the system. Shocks change user parameters and fixed variables,
and the results compare each solution with the starting levels.

The commands that work on any model work on a standalone model: macros,
attributes of variables, inequalities, the solution and equation listings,
scenario loops, the optimiser with expressions as objectives and targets,
saving and loading, GDX files and the export to GAMS, which writes the model
as a GAMS program of its own. The commands that need the built-in equations,
from \manref{cge_calibrate} to \manref{cge_report} and \manref{cge_dynamic},
stop with an error.

The example \cmd{examples/standard_model.do} writes the standard model for
the synthetic SAM in the language: the SAM is read as data, the sets of
accounts and roles from it, the parameters by assignments such as
\begin{stcode}
cge_parameter tm(c)$(sam(row, c) > 0) = sam(tax_imp, c) / sam(row, c)
cge_parameter dva(f,a) = QF0(f,a) ^ (1 + rva(a)) / sum(f2, QF0(f2,a) ^ (1 + rva(a)))
\end{stcode}
and the equations of chapter~\ref{chap:equations} that the SAM uses, 93 in
all, such as
\begin{stcode}
cge_equation T5(c): QM(c) = QD(c) * (PD(c) / PM(c) * dq(c) / dqc(c)) ^ (1 / (1 + rq(c)))
\end{stcode}
Its starting levels solve the base, and its solutions agree with those of
the built-in model to about $10^{-12}$ in four scenarios
(\chapref{chap:validation}).

\section{Nonlinear programmes}\label{sec:nlp}

From version 0.15.0, a standalone model with more free variables than
equations is a nonlinear programme when \manref{cge_solve} is given an
objective:
\begin{equation*}
\min_x\; f(x) \quad\text{subject to}\quad \ell_k(x) - r_k(x) = 0,\ k = 1,
\dots, K, \qquad L \le x \le U,
\end{equation*}
with $x$ the free variables, the fixed ones held at their levels, $f$ the
expression of \cmd{minimize()} or the negative of that of \cmd{maximize()},
and $L$ and $U$ the bounds of \manref{cge_bound},
as in GAMS's \cmd{solve m using nlp minimizing obj}. The objective is compiled
as an equation is and differentiated exactly, by forward differentiation of
its postfix rows; the equations have the exact Jacobian of
section~\ref{sec:jacobian}. An inequality is written as an equation with a
slack bounded at zero, $g(x) - S = b$ with $S \ge 0$ for $g(x) \ge b$;
inequalities paired with variables are complementarity conditions and stop
the solve with an error, as do a model without degrees of freedom and a
variable whose bounds are equal, which should be fixed.

\textbf{Algorithm.} A primal-dual interior-point method after
\citet{wachter2006}, the method of IPOPT, solves the problem in the scaled
variables $z = x/s$ and scaled residuals $c(z)$ of \manref{cge_solve}, with
the objective divided by $\max(1, \max_j |\partial f/\partial z_j| / 100)$
at the start. The start is moved inside the bounds by $\min(10^{-2}\max(1,
|b|), (U - L)/4)$ from each bound $b$; the multipliers of the equations
start at their least-squares values and those of the bounds at $\mu/(z - L)$
and $\mu/(U - z)$, with the barrier parameter $\mu = 0.1$. At each iterate
the Hessian $W$ of the Lagrangian comes from central differences of its exact
gradient with steps $10^{-6}\max(1, |z_j|)$, one-sided near a bound; the
step solves the primal-dual system
\begin{equation*}
\begin{pmatrix} W + \Sigma + \delta I & J' \\ J & 0 \end{pmatrix}
\begin{pmatrix} \Delta z \\ \lambda_+ \end{pmatrix} =
-\begin{pmatrix} \nabla f - \mu (Z - L)^{-1} e + \mu (U - Z)^{-1} e \\ c
\end{pmatrix},
\end{equation*}
with $\Sigma$ the ratios of the bound multipliers to their slacks and
$\delta$ the smallest shift, above $10^{-8}$ of the largest diagonal entry,
that makes $W + \Sigma + \delta I$ positive definite on the null space of
$J$, found from a QR decomposition of $J'$ (a small negative diagonal is
added to the lower block if the system is singular). Steps keep the
variables and bound multipliers within $\tau = \max(0.99, 1 - \mu)$ of their
bounds; the step length halves until the $\ell_1$ merit function $f - \mu
\sum \ln(\text{slacks}) + \nu \lVert c \rVert_1$, $\nu \ge \max |\lambda| +
1$, falls by Armijo's rule, a line search in place of IPOPT's filter. $\mu$
falls to $\max(\text{tol}/10, \min(0.2\mu, \mu^{1.5}))$ each time the barrier
problem is solved to $10\mu$, and the method stops when the largest of the
scaled dual residual, the residuals of the equations and the complementarity
of the bounds is below the tolerance.

\textbf{Results.} The solution is stored as a simulation with method
\cmd{nlp} and its objective, which \manref{cge_results} and the listings
report as any other. The marginals, \res{r(marginals)}, are the changes in the
optimal objective per unit increase in the right side of each equation, as in
GAMS's solution listing; the test suite checks them against differences of
the optimum in the right sides. The objective of a simulation is saved with
it, and \manref{cge_gams} exports it as \cmd{OBJV} with the equation
\cmd{EOBJV}, solved by \cmd{solve cgeprog using nlp} with IPOPT unless
\opt{}{solver()} names another NLP solver; the base, the starting levels, is
written but not solved. The export check (\texttt{tools/gams}) solves
problem 71 of \citet{hock1981}, a balancing by cross entropy and a frontier
with a binding bound in GAMS with IPOPT and compares the solutions.

\section{Integer variables}\label{sec:integer}

From version 0.18.0, \cmd{cge_bound X, integer} or \cmd{binary} (with bounds 0
and 1) makes elements of a standalone model's variables integers, and
\cmd{continuous} makes them continuous again; a nonlinear programme with such
elements is mixed-integer, as GAMS's MINLP and MIP models are. Branch and
bound solves it: each node is the relaxation with bounds on the integer
variables, solved by the interior point, or by Newton's method when the
bounds leave as many free variables as equations; a node whose relaxation is
fractional branches on the most fractional variable $x_j = v$ into $x_j \le
\lfloor v \rfloor$ and $x_j \ge \lceil v \rceil$, exploring first the child
nearer $v$, depth first; a node whose bound is no better than the incumbent,
less the gap tolerance $\max(10^{-9}, \text{optcr}) \times |f|$, is pruned.
The search ends with model status 1 (optimal) and an integer solution, whose
continuous variables are solved again with the integers fixed, or 10
(integer infeasible); a node or time limit gives 8 (integer solution) or 14
(no solution returned). The interior point finds a local optimum of each
relaxation, so that branch and bound is exact for problems whose relaxations
are convex, as linear and convex quadratic ones are, and a heuristic
otherwise. \manref{cge_gams} exports the problem as a GAMS MINLP (or, with
\cmd{type(mip)}, a MIP), declaring integer and binary variables as such,
solved by SCIP unless another solver is named.

\section{Named models}\label{sec:namedmodels}

From version 0.17.0, \cmd{cge_model define} names a model over a subset of a
standalone model's equations, as GAMS's \cmd{model} statement:
\begin{stcode}
cge_model define markets = DEM SUP
cge_model define whole = markets TOT
cge_solve, name(mk) model(markets)
\end{stcode}
A list names user equations, variables defined by an expression (their
defining equations), inequalities (their constraints) and models defined
before, or \cmd{all}. Solved on its own, the model is the system
\begin{equation*}
\ell_k(x) - r_k(x) = 0, \quad k \in \mathcal{K}_m, \qquad
x_j \text{ free for } j \in \mathcal{J}_m, \quad x_j = \bar x_j \text{ otherwise},
\end{equation*}
with $\mathcal{K}_m$ the elements of its equations, $\mathcal{J}_m$ the free
variables that appear in them, from the pattern of the Jacobian, and $\bar
x$ the levels at the start: the current levels or, with \opt{}{from()},
those of a stored simulation. Recursive and block-by-block solves, which
GAMS programs write as sequences of \cmd{solve} statements over named
models, follow: a block solved first, the next from its solution, with the
variables of the first held at their levels by \opt{}{hold()}, as GAMS's
\cmd{.fx}.
\manref{cge_model} gives an example; the simulation records its model, and
\manref{cge_gams} exports it with the equations of the model and the other
variables fixed.

\section{Error positions}\label{sec:errorpos}

From version 0.21.0, a syntax error in an expression, such as an operand or
an operator missing, two operands in a row or a parenthesis not closed, shows
the expression as it was read, its tokens separated by spaces, with a marker
under the place, before the message, as GAMS's \cmd{\$} errors do:
\begin{stcode}
. cge_equation E: X + * 2 = 1
    X + * 2
        ^
an operand is missing before *
\end{stcode}
An equation over a set is read element by element, so the marker points into
the first element at fault.

\section{Compilation}

Each equation is written out for every element of its set that meets its
condition, each sum over the elements of its set that meet its condition, and
each conditional term as itself or 0, then compiled to postfix order: numbers, variables
(positions in the stacked vector of variables), parameters, user parameters
(positions in the vector of their values), operators and functions. Stata evaluates the postfix rows at every point. From version
0.7.0 their derivatives are exact, by forward differentiation of the postfix
rows (section~\ref{sec:jacobian}): each value on the stack carries its
derivatives with respect to the variables that the element uses. Before, they
were found by grouped differences. The equations form a block \cmd{U} of the
system. They are saved with the model
and compiled again when it is calibrated again.

\textbf{GAMS.} \manref{cge_gams} writes the user parameters (\cmd{UP_NAME},
with their base values, which the scenario file changes), the user variables
(\cmd{U_NAME}, with their base levels) and one scalar equation per element (\cmd{UE1},
\cmd{UE2}, \dots) to the file \cmd{cge_user.inc}, which the GAMS model includes
before its model statement, and the results of the user variables to
\cmd{cge_user_write.inc}. GAMS then differentiates the same equations itself.
